\chapter*{Introduction Générale}
\addcontentsline{toc}{chapter}{Introduction Générale}

Dans un monde digitalisé, la transformation digitale est devenue un objectif stratégique pour les entreprises pour améliorer leur efficacité, renforcer leur compétitivité et répondre aux mutations de leur secteur. Les fiduciaires, en tant qu'acteurs clés de l'activité économique, ne font pas exception à cette évolution. Elles sont aujourd'hui confrontées à des défis majeurs liés à la gestion des prestations comptables, à la planification des ressources et à l'optimisation de leur processus métiers.

Dans ce contexte, TAMTAM International s'engage activement dans cette transformation numérique pour répondre aux différentes exigences de ces fiduciaires en leur offrant un écosystème pour satisfaire leurs besoins.

C'est dans ce cadre que s'inscrit mon projet de fin d'études, qui consiste à \textbf{Conception et développement d'un système de synchronisation incrémentale}, utilisé par le super admin pour assurer la cohérence des données et aussi la disponibilité dans notre système interne.

L'objectif essentiel de ce projet est de créer un système de synchronisation qui tourne chaque jour et aussi chaque temps qu'il exécute il donne la fraiche donnée existe sans passer par l'ancien données. Cette implémentation ne vise pas seulement de créer un nouveau système de synchronisation qui récupère les plus récents, mais également à ne pas consommer beaucoup de ressources et faire sa sans bloquer application.

\section*{Organisation du rapport}

Le présent rapport s'articule autour des chapitres suivants :

\begin{itemize}
    \item \textbf{Chapitre 1} : présente le contexte général du projet, l'organisme d'accueil et les objectifs du stage.
    
    \item \textbf{Chapitre 2} : contient l'analyse et la conception du projet.
    
    \item \textbf{Chapitre 3} : définit l'analyse des besoins fonctionnels et non fonctionnels du système ainsi que l'analyse technique qui détaille les technologies utilisées, l'architecture mise en place, ainsi que la structure de projet.
    
    \item \textbf{Chapitre 4} : illustre la solution développée à travers des interfaces représentatives.
    
    \item \textbf{Chapitre 5} : enfin on intéresse au bilan professionnel, technique et personnel.
end{itemize}
